\renewcommand{\t}{\paramValue}
\newcommand{\mdim}{d}
\newcommand{\msize}{n}
\newcommand{\avgPostElem}{\phi(\t[i])}
\newcommand{\sumDim}{\sum_{i=1}^{\mdim}}
\newcommand{\prodDim}{\prod_{i=1}^{\mdim}}
\newcommand{\vSeq}[2]{(#1_1, #1_2, \ldots, #1_{#2})}
\newcommand{\multi}{\frac{\msize!}{\prodDim \x[i]!} \prodDim \t_{i}^{\x[i]}}
\renewcommand{\H}{ \entropy{\X \mid \t}{} }
\newcommand{\omegaElem}{w}
\newcommand{\omegaElemAlt}{w^{(ij)}}
\newcommand{\eqSys}{\W(\t) = \w}
% Clique tree commands
\newcommandx{\msg}[3][1=i, 2=j, 3=\delta]{\mathFix{#3_{#1 \rightarrow #2}}}
\newcommand{\msgAlt}{\msg[i][j][\sigma]}
\newcommand{\ct}[1]{\mathFix{\mathcal{T}_{#1}}}
\newcommand{\rangeSep}{\set{r}_{i,j}}
\newcommand{\altRangeSep}{\set{r}_{i,j}^{*}}
\newcommandx{\sep}[2][1=i, 2=j]{\mathFix{\set{s}_{#1,#2}}}
\newcommand{\sepVal}{\mathFix{s_{i,j}}}
\newcommand{\msgMap}{\mathnormal{f}_{i,j}}
\newcommand{\eqScope}{\vartheta}
\newcommand{\ind}{\xi}
\newcommand{\pr}[1]{pr_{#1}}
\newcommand{\tInd}{I_{<(i, j)}}
\newcommand{\rankEqSys}{r}
\renewcommand{\M}{\psi}
\newcommand{\mat}{M_{\W}}
\newcommand{\C}[1]{C_{#1}}
\newcommandx{\wPart}[3][1=i, 2=j, 3=\set{w}]{#3_{<(#1,#2)}}
\newcommand{\wSep}{\wPart[i][j][\set{w}^*]}
\newcommand{\eqScopeSep}[2]{\eqScope_{#1,#2}}
\newcommand{\wSepVal}{\wPart[i][j][w^*]}
\newcommand{\Nb}{Nb_i}

In this section, we examine the derivation of the posterior potential for multinomial models, which are
central to our work.
By de Finetti’s theorem, any exchangeable sequence of discrete random variables can be represented using a multinomial model.
Exchangeable sequences are fundamental in statistical learning as they provide a framework for modeling data
with inherent independence and order invariance, which is common in real-world learning tasks.

A multinomial model is a parametric model (Definition~\ref{def:param-model}) in which the likelihood function follows a
multinomial distribution.
The configuration of the model is determined by two variables: $\msize$ and $\mdim$.
Each element of the sample space $s \in \set{S}$ can be represent by a $\mdim$ dimensional random vector
\[\X = \vSeq{\x}{\mdim},\] where each \x[i] is a nonnegative integer such that $\sumDim \x[i] = \msize$.
The model parameter is a nonnegative $\mdim$ dimensional real vector \[\t = \vSeq{\theta}{\mdim},\]
where $\sumDim \t[i] = 1$.
The likelihood distributing $\pi(\X \mid \theta)$ is a multinomial distribution
\begin{equation*}
    \pi(\x \mid \t) = \multi \cdot
\end{equation*}

As $\msize$ increases, the posterior distribution of \t converges to the same distribution regardless of the prior.
Therefore, when $\msize$ is sufficiently large, we can use a uniform prior to obtain the posterior potential of the
multinomial parameter.
The uniform prior is conjugate to the multinomial distribution, and the resulting posterior will follow a Dirichlet
distribution:
\begin{align*}
    p(\t \mid \x) &\sim \text{Dir}(\x[1] + 1, \x[2] + 1, \ldots, \x[\mdim] + 1) \\
    &= \prod_{j = 1}^{\mdim - 1} (\msize + j) \multi \cdot
\end{align*}
This implies that for the posterior potential we have
\begin{align*}
    \avgPost[\t] &= \avgPostExpand{\t}{\x} \\
    &= e^{\H} \prod_{j = 1}^{\mdim - 1} (\msize + j),
\end{align*}
where $\H$ is the entropy of a multinomial distribution with parameter \t.
We can write $\H$ in the following form:
\[
    \H =  - \ln(n!) + \sumDim \ln \avgPostElem,
\]
where
\begin{equation}
    \avgPostElem = \exp \left\{ -n \t[i] \ln \t[i] +
    \sum_{j = 0}^{n} \ln(j!) \binom{n}{j} \t[i]^{j} (1 - \t[i])^{n - j} \right\} \cdot
    \label{eq:pp_multi_elem}
\end{equation}
Therefore, for the posterior potential of the multinomial parameter we have
\begin{equation}
    \avgPost[\t] = (\msize + \mdim - 1)! \prodDim \avgPostElem \cdot
    \label{eq:pp_multi}
\end{equation}

For the central probability measure of a multinomial model, we usually consider a query set that consists of random
variables corresponding to fusions of the multinomial categories. %todo: should be rewrittten later

Let $\X = \vSeq{\x}{\mdim}$ be a random vector representing an element of the sample space of a multinomial model.
A fusion random variable \Y is obtained by summing up some entries of \X.
More precisely, if $\set{A}_1, \ldots, \set{A}_{\mdim'}$ is a partition of the set $\{1,2,\ldots, \mdim\}$, each value
of \Y is a vector $\vSeq{y}{\mdim'}$ where
\[
    \rval[j]{y} = \sum_{i \in \set{A}_j} \x[i], \qquad j = 1, \ldots, \mdim'\ \cdot
\]
The fusion variable \Y has a multinomial distribution with the parameter $\w = \vSeq{\omegaElem}{\mdim'}$, where
\[
    \omegaElem_j = \sum_{i \in \set{A}_j} \t[i], \qquad j = 1, \ldots, \mdim'\ \cdot
\]
For example $(\t[1] + \t[5], \t[2], \t[4] + \t[7] + \t[8],  \ldots)$ fully parametrizes the fusion variable
$(\x[1] + \x[5], \x[2], \x[4] + \x[7] + \x[8],  \ldots)$.

Suppose that in a multinomial model we aim to find the central probability measure for a query set
$\Qset = \{\X, \seqRv[Y]{k}\}$, where each \Y[i] is a fusion variable, and \X represents the sample set.
For using theorem~\ref{th:fixed_point}, we need to obtain $\avgPost$, where \w[i] is the parameter of \Y[i],
and also $\avgPost[\w]$, where $\w = \vSeq{\w}{k}$.
Note that each \w[i] is a multinomial parameter.

Each $\avgPost$ can be easily obtained using Equation~\ref{eq:pp_multi}.
However, calculating $\avgPost[\w]$ requires integration:
\[
    \avgPost[\w] = \int_{\omegaArea} \avgPost[\t] d\t \cdot
\]
This integration is performed over the set of all nonnegative solutions to the linear system of equations $\eqSys$.
Here, \W is a linear map with a matrix $\mat$ consisting solely of $0$ and $1$ entries.
The range of \W typically has much lower dimensionality than \t, and the system of equations contains
many free variables.
The equation system has a solution when \w is inside the convex hull of the columns of $\mat$.

We convert this integration into an inference problem in a probabilistic graphical model, and then use a slightly
modified version of the clique tree message passing~\cite{koller2009probabilistic} algorithm.

We represent each equation in the linear system $\eqSys$ using an indicator factor.
Let $\eqScope \subseteq \vSeq{\t}{\mdim}$, and $0 \leq \omegaElem \leq 1$.
An indicator factor $\ind$ is a function with the scope $\{\eqScope, \omegaElem\}$ such that:
\[
    \ind(\eqScope, \omegaElem) = \begin{cases}
                                     1 \quad &  \sum_{\t[i] \in \eqScope} \t[i] = \omegaElem, \\
                                     0 & \text{otherwise} \cdot
    \end{cases}
\]
For example $\ind(\t[2], \t[5], \t[6], \omegaElem)$ is $1$ when $\t[2] + \t[5] + \t[6] = w$, and is $0$ otherwise.

Suppose the rank of $\mat$ is $\rankEqSys$.
We select $\rankEqSys$ linearly independent equations from $\eqSys$, and define an indicator factor for every
equation.
Then, we construct an unnormalized measure $\M$:
\[
    \M(\t, \omegaElem_1, \ldots, \omegaElem_{\rankEqSys}) =
    \prodDim \avgPostElem \prod_{j = 1}^{\rankEqSys} \ind(\eqScope_j, \omegaElem_j) \cdot
\]
Clearly, marginalizing $\M$ over \t yields $\avgPost[\w]$.
To achieve this, we construct a clique tree for $\M$ that has at least one clique $\C{\omegaElem}$ containing all
$\omegaElem_j$.
By performing message passing and selecting $\C{\omegaElem}$ as the root of the tree, we can compute the desired
marginal: $\M(\omegaElem_1, \ldots, \omegaElem_{\rankEqSys})$.
We denote this clique tree by $\ct{\M}$.

In $\ct{\M}$ each equation and each variable of the equation system $\eqSys$ is associated with a clique.
Let $\C{i}$ and $\C{j}$ be two adjacent cliques in $\ct{\M}$ with the sepset \sep.
Without loss of generality, we assume that the root is in the $\C{j}$ side of the tree.
Let $\wPart$ be the set of all $\omegaElem_k$ such that an indicator $\ind(\cdot, \omegaElem_k)$ is associated
with a clique in the $\C{i}$ side of the tree.
We can think of $\wPart$ as the set of all equations that are associated with the $\C{i}$ side of the tree.
Let $\eqScopeSep{i}{j}$ be the set of all \t[i], i.e., the variables of the equation system that are in the scope of
both $\C{i}$ and $\C{j}$.
Clearly we have \[ \sep = \wPart \cup \eqScopeSep{i}{j} \cdot \]
However, it turns out that during message passing, we can use a simpler sepset.

\begin{definition}
    Let $\ct{}$ be a clique tree.
    Let $\rangeSep$ denote the set of all possible assignments to the sepset \sep.
    A \emph{message reparameterization} for $\ct{}$ is a mapping $\msgMap$ defined from $\rangeSep$
    to another set $\altRangeSep$ for every $i, j$, which is message preserving.
    That means, for any $\sepVal, \sepVal' \in \rangeSep$, if $\msgMap(\sepVal) = \msgMap(\sepVal')$
    then $\msg(\sepVal) = \msg(\sepVal')$.

    The reparameterization of a message \msg is a function \msgAlt defined on $\altRangeSep$ such that for every
    $\sepVal \in \rangeSep$, we have $\msg(\sepVal) = \msgAlt(\msgMap(\sepVal))$.
\end{definition}

\begin{proposition}
    We can replace all messages \msg with their reparameterization \msgAlt without affecting the message passing
    algorithm.
\end{proposition}

\begin{proof}
    During message passing no information regarding specific assignments to \sep is passed between cliques.
    $\C{i}$ simply produces \msg for every possible assignments to \sep.
    Having the reparameterization mapping the reparameterized messages can be used to produce \msg.
\end{proof}

Suppose $\omegaElem \in \wPart$ and $\ind(\eqScope, \omegaElem)$ is the indicator containing
$\omegaElem$.
This indicator represents the equation
\[
    \sum_{\t[i] \in \eqScope} \t[i] = \omegaElem \cdot
\]
We define a new variable $0 \leq \omegaElemAlt \leq 1$ for every $\omegaElem \in \wPart$ such that
\[
    \omegaElemAlt = \omegaElem - \sum_{\t[i] \in \eqScope \cap \eqScopeSep{i}{j}} \t[i]
\]
For each sepset \sep, there will be a corresponding set of these new variables.
We denote this set by $\wSep$.
We have a map from assignments to \sep into assignments to $\wSep$, which is message preserving.
If two assignments $\sepVal$, $\sepVal'$ to the sepset variables induce the same assignment to
$\wSep$ variables, that will imply that the equation system corresponding to the $\C{i}$ part of the tree has
the same solution set for $\sepVal$ and $\sepVal'$.
As a result $\msg(\sepVal) = \msg(\sepVal')$.

To simplify message computation, we may impose some restrictions on the clique tree structure.
For each clique $\C{i}$, Let $\pr{i}$ be the upstream neighbor of $i$, the one on the path to the root clique
$\C{\omegaElem}$.
We require that $\C{i} - \sep[i][\pr{i}] = \{\t[i]\}$.
That is, in each clique exactly one variable should be eliminated.
We associate $\avgPostElem$ with the clique that eliminates \t[i].
Let $\Nb$ be the set of indexes of cliques that are neighbors of $\C{i}$.
We define $\tInd$ to be an indicator vector.
The message from $\C{i}$ to another clique $\C{j}$ is computed using the following convolution:
\begin{equation}
    \msg(\wSepVal) = \int \avgPostElem \prod_{k \in \Nb - \{j\}} \msg[k][i](\wSepVal - \t[i] \tInd) d\t[i]
\end{equation}
As we discussed in the previous section, we usually need to compute the
posterior potential of the joint feature $\W(\paramValue) = \vSeq{\w}{m}$, where each \w[i] is a function of the model
parameter \paramValue that fully parametrizes some random variable \X[i] in the model.

For example, $(\t[1] + \t[4], \t[2] + \t[3] + \t[5],  \ldots)$ fully parametrizes $(\x[1] + \x[4], \x[2] + \x[3] + \x[5],  \ldots)$.
When we consider fusion variables, $\W(\t)$ will be a linear map that only includes addition.
The matrix of this mapping will include only $0$ and $1$.

We aim to compute the posterior potential of the joint features \W

is a linear map of the model
parameter that only
includes addition.
We can represent this linear map with a matrix, this matrix will only include 1 and 0
We aim to compute the posterior potential of the joint features \W in a multinomial model.
We can think of $\eqScopeSep{i}{j}$ as the set of all variables of the equation system that are associated with
the $\C{i}$ side of the tree.

We may represent \w with a column vector
\[
    \w = \begin{bmatrix}
             \W[1](\paramValue) \\
             \W[2](\paramValue) \\
             \vdots             \\
             \W[m](\paramValue) \\
    \end{bmatrix} =
    \begin{bmatrix}
        \w[1]  \\
        \w[2]  \\
        \vdots \\
        \w[m]  \\
    \end{bmatrix}
\]


\[
    \avgPost[\t] \propto A e^B \prod_{i=1}^{k} \avgPostElem
\]


\[
    \avgPost[\t] \propto A e^{\entropy{\X \mid \t}{}}
\]

\[
    \entropy{\X \mid \t}{} = B + \sum_{i=1}^{k} \ln \avgPostElem
\]
\[
    B = - \ln(n!)
\]
\[
    \avgPost[\t] \propto \prod_{i=1}^{k} \avgPostElem
\]
